% Shared preamble of the ArchytasZX reference manual.
\usepackage[T1]{fontenc}
\usepackage[utf8]{inputenc}
\usepackage{lmodern}
\usepackage{microtype}
\usepackage[a4paper,margin=2.5cm]{geometry}
\usepackage{amsmath,amssymb,amsthm,mathtools,stmaryrd}
\usepackage{graphicx}
\usepackage{booktabs,tabularx,longtable,xltabular,array,multirow}
\usepackage{xcolor}
\usepackage{listings}
\usepackage[most]{tcolorbox}
\usepackage{float}
\usepackage{enumitem}
\usepackage{csquotes}
\usepackage{tikz}
\usetikzlibrary{shapes.geometric,shapes.misc,positioning,calc,fit,backgrounds,arrows.meta,decorations.pathreplacing}
\usepackage[numbers,sort&compress]{natbib}
\usepackage{imakeidx}
\makeindex[intoc,columns=2]
\usepackage[hidelinks,pdfusetitle,bookmarksnumbered]{hyperref}
\usepackage[nameinlink,capitalise,noabbrev]{cleveref}

\setlist{itemsep=2pt,topsep=4pt}
\setcounter{secnumdepth}{3}
\setcounter{tocdepth}{2}
\raggedbottom
\setlength{\emergencystretch}{3em}

% Theorem-like environments.
\theoremstyle{definition}
\newtheorem{definition}{Definition}[chapter]
\newtheorem{algorithm}[definition]{Algorithm}
\newtheorem{example}[definition]{Example}
\theoremstyle{plain}
\newtheorem{proposition}[definition]{Proposition}
\newtheorem{theorem}[definition]{Theorem}
\newtheorem{lemma}[definition]{Lemma}
\theoremstyle{remark}
\newtheorem{remark}[definition]{Remark}
\crefname{algorithm}{Algorithm}{Algorithms}

% Mathematical notation, as in the paper.
\newcommand{\ket}[1]{\lvert #1 \rangle}
\newcommand{\bra}[1]{\langle #1 \rvert}
\newcommand{\braket}[2]{\langle #1 \vert #2 \rangle}
\newcommand{\sem}[1]{\llbracket #1 \rrbracket}
\newcommand{\C}{\mathbb{C}}
\newcommand{\N}{\mathbb{N}}
\newcommand{\Z}{\mathbb{Z}}
\newcommand{\Q}{\mathbb{Q}}
\newcommand{\R}{\mathbb{R}}
\newcommand{\Zd}[1]{\mathbb{Z}_{#1}}
\newcommand{\bbox}{!-box}
\newcommand{\bboxes}{!-boxes}

% Code listings.
\definecolor{azxcmd}{RGB}{0,70,140}
\definecolor{azxgen}{RGB}{0,115,70}
\definecolor{azxkw}{RGB}{125,40,110}
\definecolor{azxcomment}{gray}{0.45}
\definecolor{leankw}{RGB}{125,40,110}
\lstset{
  basicstyle=\ttfamily\small,
  columns=fullflexible,
  keepspaces=true,
  frame=single,
  framerule=0.4pt,
  xleftmargin=1em,
  xrightmargin=1em,
  captionpos=t,
  breaklines=true,
  upquote=true,
  literate={→}{{$\to$}}1 {←}{{$\leftarrow$}}1 {≤}{{$\le$}}1 {≥}{{$\ge$}}1 {≠}{{$\ne$}}1
           {ω}{{$\omega$}}1 {⟩}{{$\rangle$}}1 {⟨}{{$\langle$}}1 {∀}{{$\forall$}}1
           {∃}{{$\exists$}}1 {λ}{{$\lambda$}}1 {∑}{{$\sum$}}1 {·}{{$\cdot$}}1
           {ℂ}{{$\mathbb{C}$}}1 {ℕ}{{$\mathbb{N}$}}1 {ℤ}{{$\mathbb{Z}$}}1 {ℚ}{{$\mathbb{Q}$}}1
           {↦}{{$\mapsto$}}1 {⊗}{{$\otimes$}}1 {∘}{{$\circ$}}1 {≫}{{$\gg$}}1
           {α}{{$\alpha$}}1 {β}{{$\beta$}}1 {π}{{$\pi$}}1 {∈}{{$\in$}}1 {×}{{$\times$}}1
           {₀}{{$_0$}}1 {₁}{{$_1$}}1 {₂}{{$_2$}}1 {⁻}{{$^{-}$}}1 {¹}{{$^{1}$}}1
           {√}{{$\surd$}}1 {∧}{{$\wedge$}}1 {∨}{{$\vee$}}1 {¬}{{$\neg$}}1 {↔}{{$\leftrightarrow$}}1
           {⟶}{{$\longrightarrow$}}1 {ε}{{$\varepsilon$}}1 {μ}{{$\mu$}}1 {σ}{{$\sigma$}}1
           {ℝ}{{$\mathbb{R}$}}1 {θ}{{$\theta$}}1 {τ}{{$\tau$}}1 {ρ}{{$\rho$}}1 {ν}{{$\nu$}}1
           {ζ}{{$\zeta$}}1 {φ}{{$\varphi$}}1 {ψ}{{$\psi$}}1 {γ}{{$\gamma$}}1 {δ}{{$\delta$}}1
           {Δ}{{$\Delta$}}1 {Φ}{{$\Phi$}}1 {Ψ}{{$\Psi$}}1 {Σ}{{$\Sigma$}}1 {Π}{{$\Pi$}}1
           {∏}{{$\prod$}}1 {−}{{$-$}}1 {∉}{{$\notin$}}1 {∣}{{$\mid$}}1 {⊆}{{$\subseteq$}}1
           {≡}{{$\equiv$}}1 {…}{{$\ldots$}}1 {₃}{{$_3$}}1,
}
\lstdefinelanguage{azx}{
  sensitive=true,
  morekeywords=[1]{load,show,rules,match,apply,check,normalize,saturate,decide,prove,induct,
    certificate,translate,name,names,undo,reset,detail,help,double,purify},
  morekeywords=[2]{dim,count,let,wire,inputs,outputs,scalar,box,repeat,as,in,phase,weight,
    literal,against,for,on,exhaustive,algebraic,emit,verify,kernel,lean,back},
  morekeywords=[3]{Z,X,W,F,T,Ti,S,B,G},
  morecomment=[l]{\#},
  morestring=[b]",
}
\lstdefinestyle{azx}{
  language=azx,
  keywordstyle=[1]{\color{azxcmd}},
  keywordstyle=[2]{\color{azxkw}},
  keywordstyle=[3]{\color{azxgen}},
  commentstyle={\itshape\color{azxcomment}},
}
\lstdefinelanguage{Lean}{
  morekeywords={theorem,lemma,def,abbrev,structure,inductive,class,instance,by,fun,let,in,where,
    match,with,axiom,import,namespace,section,end,open,variable,noncomputable,private,
    protected,example,deriving,if,then,else,do,have,show,from,calc,decide,simp,rfl},
  sensitive=true, morecomment=[l]{--}, morecomment=[s]{/-}{-/},
}
\lstdefinestyle{lean}{
  language=Lean,
  keywordstyle={\color{leankw}},
  commentstyle={\itshape\color{azxcomment}},
}
\lstdefinestyle{python}{
  language=Python,
  keywordstyle={\color{azxcmd}},
  commentstyle={\itshape\color{azxcomment}},
}
\lstdefinestyle{shell}{language=bash, keywordstyle={}}
\lstdefinestyle{grammar}{language={}, basicstyle=\ttfamily\small}
\lstdefinestyle{inline}{frame=none, basicstyle=\ttfamily, keywordstyle={}, breaklines=true,
  breakatwhitespace=false}

% Inline code: \azx!...! (engine language), \py!...! (Python), \lean!...! (Lean),
% \code!...! (other literal text), \file{...} (paths), \module{...} (Python modules).
\newcommand{\azx}{\lstinline[style=azx,frame=none,basicstyle=\ttfamily,breaklines=true,breakatwhitespace=false]}
\newcommand{\nat}{\azx}
\newcommand{\py}{\lstinline[style=inline,language=Python]}
\newcommand{\code}{\lstinline[style=inline]}
\newcommand{\lean}{\lstinline[style=lean,frame=none,basicstyle=\ttfamily,breaklines=true,breakatwhitespace=false]}
\def\UrlBreaks{\do\/\do\.\do\_\do\-}
\newcommand{\file}[1]{\nolinkurl{#1}}
\newcommand{\modidx}[1]{\nolinkurl{#1}}
\newcommand{\module}[1]{\nolinkurl{#1}\index{#1@\modidx{#1}}}
\newcommand{\term}[1]{\emph{#1}\index{#1}}
\newcommand{\tool}[2]{#1\footnote{\url{#2}}}

% Specification boxes for signatures and record layouts.
\newtcolorbox{signature}[1][]{colback=gray!4,colframe=gray!50,boxrule=0.4pt,arc=1pt,
  left=4pt,right=4pt,top=2pt,bottom=2pt,fontupper=\small,#1}

% ZXW diagram styles. Diagrams flow bottom to top: inputs at the bottom, outputs at the top.
\definecolor{zxgreen}{RGB}{204,255,204}
\definecolor{zxred}{RGB}{255,136,136}
\definecolor{zxyellow}{RGB}{255,255,150}
\tikzset{
  zxd/.style={baseline=(current bounding box.center),every node/.style={font=\scriptsize},
    x=0.9cm,y=0.9cm},
  wire/.style={draw,line width=0.5pt},
  zs/.style={circle,draw,fill=zxgreen,minimum size=4.5mm,inner sep=0.5pt,font=\scriptsize},
  xs/.style={circle,draw,fill=zxred,minimum size=4.5mm,inner sep=0.5pt,font=\scriptsize},
  zdot/.style={circle,draw,fill=zxgreen,minimum size=2.2mm,inner sep=0pt},
  xdot/.style={circle,draw,fill=zxred,minimum size=2.2mm,inner sep=0pt},
  fbox/.style={rectangle,draw,fill=zxyellow,minimum size=4mm,inner sep=1pt,font=\scriptsize},
  tri/.style={isosceles triangle,isosceles triangle apex angle=60,shape border rotate=90,
    draw,fill=white,minimum size=3.5mm,inner sep=0pt},
  trii/.style={isosceles triangle,isosceles triangle apex angle=60,shape border rotate=-90,
    draw,fill=white,minimum size=3.5mm,inner sep=0pt},
  wnode/.style={isosceles triangle,isosceles triangle apex angle=60,shape border rotate=90,
    draw,fill=black,minimum size=3mm,inner sep=0pt},
  bnode/.style={trapezium,trapezium angle=65,draw,fill=gray!20,minimum width=7mm,
    minimum height=3.5mm,inner sep=1pt,font=\scriptsize},
  snode/.style={trapezium,trapezium angle=115,draw,fill=gray!20,minimum width=7mm,
    minimum height=3.5mm,inner sep=1pt,font=\scriptsize},
  ground/.style={draw,line width=0.8pt},
  bang/.style={draw,dashed,rounded corners=2pt,inner sep=3pt},
  banglabel/.style={font=\tiny,anchor=north west,inner sep=1pt},
  dimlabel/.style={font=\tiny,inner sep=1pt},
}
% \zxground{coordinate}: the discard G drawn as an earth symbol above the given point.
\newcommand{\zxground}[1]{%
  \draw[ground] ($(#1)+(-0.18,0)$) -- ($(#1)+(0.18,0)$);
  \draw[ground] ($(#1)+(-0.11,0.07)$) -- ($(#1)+(0.11,0.07)$);
  \draw[ground] ($(#1)+(-0.04,0.14)$) -- ($(#1)+(0.04,0.14)$);}
